\subsection{相对有界线性形式}

设 E 是有向序向量空间。设 Q 是 E 上正线性形式所成的集；它是 E 的代数对偶 \(E^{*}\) （即 E 上一切线性形式所成的空间）的一个子集。直接可知 \(Q + Q \subset Q\) 且 \(\lambda Q \subset Q\) 对每个标量 \(\lambda > 0\) 成立（换句话说，Q 是 \(E^{*}\) 中的一个凸锥）。此外，\(Q \cap (-Q) = \{0\}\) ，因为若 L 与 -L 都是正线性形式，则 \(L(x) \geqslant 0\) 与 \(L(x) \leqslant 0\) 对一切 \(x \geqslant 0\) 成立，从而 \(L(x) = 0\) 对一切 \(x \geqslant 0\) 成立，因此 L = 0 (No. 1, Prop. 3)。于是集 Q 在 \(E^{*}\) 上定义了一个序关系 \(L \leqslant M\) ，它等价于«M - L 是 E 上的正线性形式»，或者等价于«对每个 \(x \geqslant 0\) ，\(L(x) \leqslant M(x)\) »；对于这个序结构，元素 \(\geqslant 0\) 的 \(E^{*}\) 中的成员就是正线性形式（这证明了所引进术语的合理性）。设 \(\Omega\) 是由 Q 生成的 \(E^{*}\) 的线性子空间，即 E 上作为两个正线性形式之差的线性形式所成的集；当 E 是 Riesz 空间时，我们将给出 \(\Omega\) 的元素的另一个刻划。

定义 2. --- 给定一个 Riesz 空间 E，若对 E 中每个 \(x \geqslant 0\) ，L 在由 \(y \in E\) 中满足 \(|y| \leqslant x\) 者所组成的集上有界，则称 E 上的线性形式 L 为相对有界的。

定理 1. --- \(1^{\circ}\) 为使 Riesz 空间 E 上的线性形式 L 是相对有界的，必要且充分的是：它是两个正线性形式之差。

\(2^{\circ}\) E 上相对有界线性形式所成的序向量空间 \(\Omega\) 是一个完全格序的 Riesz 空间。

若 \(L = U - V\) ，其中 \(U\) 与 \(V\) 是 \(\mathbf{E}\) 上的正线性形式，则关系 \(-x \leqslant y \leqslant x\) 蕴涵

\[
- U (x) \leqslant U (y) \leqslant U (x) \quad \text { and } \quad - V (x) \leqslant V (y) \leqslant V (x),
\]

由此立即可得 \(|L(y)| \leqslant U(x) + V(x)\) ；于是 L 是相对有界的。反过来，设 L 是相对有界的；问题归结为证明存在一个正线性形式 N，使 \(N(x) \geqslant L(x)\) 对一切 \(x \geqslant 0\) 成立，因为这时 N - L 将是一个正线性形式。

而若一个正线性形式 N 具有这一性质，则对每个 \(x \geqslant 0\) 及 \(0 \leqslant y \leqslant x\) ，我们有 \(N(x) \geqslant N(y) \geqslant L(y)\) ，因此 \(N(x) \geqslant \sup_{0 \leqslant y \leqslant x} L(y)\) ；若我们能证明实值函数

\[
x \mapsto M (x) = \sup _ {0 \leqslant y \leqslant x} L (y),
\]

（它定义在 E 中元素 \(\geqslant0\) 所成的集 P 上）可以延拓为 E 上的一个正线性形式（仍记作 M），则我们就证明了定理的第一部分，并且还证明了 M 是 \(\Omega\) 中 0 与 L 的上确界。由于在 P 上 \(M(x)\geqslant0\) ，问题归结为证明

\[
M (x + x ^ {\prime}) = M (x) + M (x ^ {\prime})
\]

对 E 的每对元素 \(x \geqslant 0\) 、\(x' \geqslant 0\) 成立 (No. 1, Prop. 3)。由定义，

\[
\begin{array}{c} M (x) + M (x ^ {\prime}) = \sup _ {0 \leqslant y \leqslant x} L (y) + \sup _ {0 \leqslant y ^ {\prime} \leqslant x ^ {\prime}} L (y ^ {\prime}) \\ = \sup _ {0 \leqslant y \leqslant x, 0 \leqslant y ^ {\prime} \leqslant x ^ {\prime}} L (y + y ^ {\prime}) \leqslant M (x + x ^ {\prime}). \end{array}
\]

另一方面，对每个 \(z\) （它满足 \(0 \leqslant z \leqslant x + x'\) ），我们有 \(x + x' = z + u\) ，其中 \(u \geqslant 0\) ；由分解引理 (§1, No. 1)，存在两个元素 \(y, y'\) ，使得 \(0 \leqslant y \leqslant x\) 、\(0 \leqslant y' \leqslant x'\) ，且 \(z = y + y'\) ，\(u = (x - y) + (x' - y')\) ；于是

\[
L (z) = L (y) + L \left(y ^ {\prime}\right) \leqslant M (x) + M \left(x ^ {\prime}\right),
\]

因此 \(M(x + x') = \sup_{0 \leqslant z \leqslant x + x'} L(z) \leqslant M(x) + M(x')\) ，这就完成了定理第一部分的证明。此外，我们这样还证明了 \(\Omega\) 是一个 Riesz 空间，并且对 E 上每个相对有界线性形式 L 及每个 \(x \geqslant 0\) ，

\[
L ^ {+} (x) = \sup _ {0 \leqslant y \leqslant x} L (y).\tag{1}
\]

剩下的是要看出 \(\Omega\) 是完全格序的；为此，只需证明正线性形式的任何有上界且依关系 \(\leqslant\) 有向的集 H 在 \(\Omega\) 中有一个上确界。

更一般地，我们有下述引理：

引理. --- 设 E 是有向序向量空间，\(E^{*}\) 是它的对偶，以正线性形式为正元素而排序。设 \((u_{\alpha})\) 是 \(E^{*}\) 中元素的一个递增有向族。若对 E 中每个 \(x \geqslant 0\) ，\(\sup u_{\alpha}(x) < +\infty\) ，则族 \((u_{\alpha})\) 在 \(E^{*}\) 中有一个上确界 u，且对 E 中一切 \(x \geqslant 0\) ，

\[
u (x) = \sup _ {\alpha} u _ {\alpha} (x).\tag{2}
\]

在 E 中一切 \(x \geqslant 0\) 所成的集 P 上，用公式 (2) 定义映射 u；直接可知，\(u(\lambda x) = \lambda u(x)\) 对一切 \(\lambda \geqslant 0\) 及 \(x \in P\) 成立；因此，由 No. 1 的命题 2，为证明本引理，只需证明

\[
u (x + y) = u (x) + u (y)
\]

对 P 中的 \(x, y\) 成立。但只要注意到，关于指标的定向集，\(u(x) = \lim u_{\alpha}(x)\) （单调极限定理），这就立即可得。

由公式 (1) 立即推出：若 L 与 M 是 E 上两个相对有界线性形式，则对每个 \(x \geqslant 0\) ，

\[
\left\{ \begin{array}{l} \sup (L, M) (x) = \sup _ {y \geqslant 0,   z \geqslant 0,   y + z = x} \big (L (y) + M (z) \big) \\ \inf (L, M) (x) = \inf _ {y \geqslant 0,   z \geqslant 0,   y + z = x} \big (L (y) + M (z) \big). \end{array} \right.\tag{3}
\]

特别地，若在这些公式的第一个中把 \(M\) 换成 \(-L\) ，就得到

\[
| L | (x) = \sup _ {y \geqslant 0, z \geqslant 0, y + z = x} L (y - z).
\]

而若 \(x = y + z\) 、\(y \geqslant 0\) 且 \(z \geqslant 0\) ，则 \(-x \leqslant y - z \leqslant x\) ；反过来，关系 \(|u| \leqslant x\) 蕴涵 \(L(u) \leqslant |L|(|u|) \leqslant |L|(x)\) 。由此我们导出公式

\[
| L | (x) = \sup _ {| y | \leqslant x} L (y) \quad \text { for } x \geqslant 0,\tag{4}
\]

从而特别有，

\[
| L (x) | \leqslant | L | (| x |)\tag{5}
\]

对一切 \(x \in \mathbf{E}\) 成立。

命题 4. --- 为使 Riesz 空间 E 上两个正线性形式 L, M 在空间 \(\Omega\) 中互奇异，必要且充分的是：对每个数 \(\varepsilon > 0\) 及 E 中每个 \(x \geqslant 0\) ，存在 E 的两个元素 \(y \geqslant 0\) 、\(z \geqslant 0\) ，使得 \(x = y + z\) 且 \(L(y) + M(z) \leqslant \varepsilon\) 。事实上，由 (3) 中第二个公式，这一条件表示 \(\inf(L, M) = 0\) 。

命题 5. --- 设 L 是 Riesz 空间 E 上的一个正线性形式。为使 E 上的正线性形式 M 属于 \(\Omega\) 中由 L 生成的带，必要且充分的是：对 E 中每个 \(x \geqslant 0\) 及每个数 \(\varepsilon > 0\) ，存在一个数 \(\delta > 0\) ，使得关系 \(0 \leqslant y \leqslant x\) 与 \(L(y) \leqslant \delta\) 蕴涵 \(M(y) \leqslant \varepsilon\) 。

我们先证明条件是必要的。若 \(M \geqslant 0\) 属于 \(\Omega\) 中由 L 生成的带，则 (§1, No. 5, Prop. 6 的推论)

\[
M = \sup _ {n} \left(\inf (n L, M)\right).
\]

若令

\[
U _ {n} = M - \inf (n L, M),
\]

于是 \(U_{n}\) 是 E 上的一个正线性形式，且 \(\inf_{n} U_{n} = 0\) 在 \(\Omega\) 中成立；因而（引理）当 n 趋于无穷时 \(U_{n}(x)\) 趋于 0，且存在 n 使 \(U_{n}(x) \leqslant \varepsilon / 2\) 。固定这样一个 n，则 \(U_{n}(y) \leqslant \varepsilon / 2\) 对一切满足 \(0 \leqslant y \leqslant x\) 的 y 成立，于是关系 \(0 \leqslant y \leqslant x\) 蕴涵

\[
M (y) \leqslant \frac {\varepsilon}{2} + \inf (n L, M) (y) \leqslant \frac {\varepsilon}{2} + n L (y);
\]

若 \(y\) 使 \(L(y) \leqslant \varepsilon / 2n\) ，则得 \(M(y) \leqslant \varepsilon\) ，这就证明了我们的断言。

现在证明条件是充分的。设 \(M\) 是 \(\mathrm{E}\) 上任一正线性形式，则可写 \(M = U + V\) ，其中 \(U\) 属于由 \(L\) 在 \(\Omega\) 中生成的带，而 \(V\) 与 \(L\) 互奇异，\(U\) 与 \(V\) 都是正的 (§1, No. 5, Th. 1)。若 \(M\) 满足命题中的条件，则 \(V = M - U\) 也满足，因为 \(0 \leqslant V \leqslant M\) 。由此我们将推出 \(V = 0\) 。对每个 \(x \geqslant 0\) （在 \(\mathrm{E}\) 中）及每个数 \(\eta > 0\) ，存在两个元素 \(y \geqslant 0\) 、\(z \geqslant 0\) （属于 \(\mathrm{E}\)），使得 \(x = y + z\) 且 \(L(y) + V(z) \leqslant \eta\) （命题 4）；给定任意数 \(\varepsilon > 0\) ，选取 \(\eta \leqslant \varepsilon\) ，使得关系 \(0 \leqslant u \leqslant x\) 与 \(L(u) \leqslant \eta\) 蕴涵 \(V(u) \leqslant \varepsilon\) ；把 \(y\) 与 \(z\) 取定如上，我们有 \(L(y) \leqslant \eta\) ，因此 \(V(y) \leqslant \varepsilon\) ，于是

\[
V (x) = V (y) + V (z) \leqslant \varepsilon + \eta \leqslant 2 \varepsilon ;
\]

由于 \(\varepsilon\) 是任意的，我们有 \(V(x) = 0\) 对每个 \(x \geqslant 0\) 成立，即 \(V = 0\) 。

例. --- 设 E 是一个 Riesz 空间，赋以一个与其序向量空间结构相容的局部凸拓扑 (TVS, II, §2, No. 7)。设 \(E'\) 是 E 的拓扑对偶，并再设 E 中元素 \(\geqslant 0\) 所成的锥 P 对于弱化拓扑 \(\sigma(\mathrm{E},\mathrm{E}')\) 是完备的。则每个连续线性形式 \(x' \in E'\) 都是相对有界的，因为已知 (TVS, II, §6, No. 8, Prop. 11 的推论 2)，在这些条件下，对 E 中每个 \(x \geqslant 0\) ，由 \(y \in E\) 中满足 \(|y| \leqslant x\) 者所成的集对 \(\sigma(\mathrm{E},\mathrm{E}')\) 是紧的。由此我们推出 E 这时是完全格序的；因为 (§1, No. 3, Prop. 2)，只需证明：对每个集 \(H \subset E\) （它有上界且依 \(\leqslant\) 有向），H 的截断滤子 F 在 E 中对拓扑 \(\sigma(\mathrm{E},\mathrm{E}')\) 收敛（后者与 E 的序向量空间结构相容）。由平移，我们可以假设 \(H \subset P\) ，于是只需证明 F 是对 \(\sigma(\mathrm{E},\mathrm{E}')\) 的 Cauchy 滤子，或者再说一遍，每个连续线性形式 \(x' \in E'\) 关于 F 有极限。而当 \(x'\) 是正线性形式时，这由单调极限定理立即可得；又由于每个线性形式 \(x' \in E'\) 是两个正线性形式之差 (Th. 1)，我们的断言得证。
% semantic-fragments: legacy-input-seam
\relax{}
